\documentclass[aspectratio=169]{beamer}

\usepackage[italian]{babel}
\usepackage[utf8]{inputenc}
\usepackage[T1]{fontenc}
\usepackage{lmodern} 
\usepackage{amsmath}
\usepackage{tikz}
\usepackage{silence}
\WarningFilter{latexfont}{Font shape}
\WarningFilter{latexfont}{Some font shapes}
\usepackage{clrscode3e}
\usetikzlibrary{shapes.geometric, arrows.meta, positioning}
\usefonttheme[onlymath]{serif}

%Colori
\definecolor{darkgreen}{RGB}{10, 134, 54}
\definecolor{lightgreen}{RGB}{13, 179, 77}
\definecolor{blockbodybg}{RGB}{232, 242, 234}

%Impostazioni Tema Beamer
\usefonttheme{default}
\setbeamertemplate{navigation symbols}{}

\setbeamercolor{structure}{fg=darkgreen}
\setbeamercolor{title}{bg=lightgreen, fg=black}
\setbeamercolor{frametitle}{bg=lightgreen, fg=white}
\setbeamercolor{block title}{bg=darkgreen, fg=white}
\setbeamercolor{block body}{bg=blockbodybg, fg=black}
\setbeamercolor{itemize item}{fg=darkgreen}
\setbeamercolor{itemize subitem}{fg=darkgreen}

\setbeamertemplate{itemize items}[ball]

\setbeamertemplate{blocks}[rounded][shadow=true]

\setbeamertemplate{title page}{
    \vfill
    \begin{center}
        \begin{beamercolorbox}[sep=8pt,center,shadow=true,rounded=true,wd=0.85\paperwidth]{title}
            \usebeamerfont{title}\inserttitle\par%
        \end{beamercolorbox}
        \vskip2.5em\par
        \usebeamerfont{institute}\insertinstitute\par
        \vskip1.5em\par
        \usebeamerfont{date}\insertdate\par
    \end{center}
    \vfill
}

% Barra superiore (Headline)
\setbeamertemplate{headline}{
    \begin{beamercolorbox}[ht=3.5ex,dp=1ex]{block title}
    \end{beamercolorbox}
}

% Barra inferiore (Footline)
\setbeamertemplate{footline}{
    \begin{beamercolorbox}[ht=2.5ex,dp=1ex]{block title}
    \end{beamercolorbox}
}


% Informazioni Titolo
\title{Breaking the Sorting Barrier for Directed Single-Source Shortest Paths}
\institute{Raed Abbas\\ Sotto la supervisione del Prof. Alessandro Barenghi\\Politecnico di Milano}
\date{5 marzo 2026}

\begin{document}

% Slide: Titolo 
\begin{frame}
    \titlepage
\end{frame}

% Slide : Introduzione
\begin{frame}{Introduzione}
    \begin{block}{Ingredienti}
        \begin{itemize}
            \item Grafo Diretto $\mathbf{G = (V,E)}$ dove:
            \begin{itemize}
                \item $\mathbf{V}$: insieme dei vertici (con $\vert \mathbf{V}\vert$=n)
                \item $\mathbf{E}$: insieme dei archi (con $\vert \mathbf{E}\vert$=m)
            \end{itemize}
            \item Lavoriamo con archi a pesi non negativi, quindi definiamo la funzione peso: \\ $w_{uv}$ $:\mathbf{E}$ $\rightarrow$ $\mathbb{R}_{\ge 0}$  $\forall (u,v) \in \mathbf{E}$
            \item Vertice sorgente $s$
        \end{itemize}
    \end{block}
    \begin{block}{Obiettivo}
    Trovare la lunghezza del percorso più breve tra $s$ e ogni vertice in $\mathbf{V}$.
    \end{block}
\end{frame}
% Slide: Dijkstra Classico
\begin{frame}[fragile]{L'approccio storico: \textsc{Dijkstra} Classico}
    \begin{columns}
        \begin{column}{0.5\textwidth}
            La versione base seleziona il minimo scorrendo l'intero vettore $\mathbf{d}$ ad ogni iterazione.
            \vspace{0.2cm}
            \begin{block}{Complessità}
                \begin{itemize}
                    \item \textbf{Estrazione minimo:} $\mathcal{O}(\vert\mathbf{V}\vert)$
                    \item \textbf{Iterazioni:} $\vert\mathbf{V}\vert$
                    \item \textbf{Complessità Computazionale:} $\mathcal{O}(\vert V\vert^2 + \vert E\vert)$
                \end{itemize}
            \end{block}
        \end{column}
        \begin{column}{0.5\textwidth}
            \scriptsize
            \begin{codebox}
            \Procname{$\proc{Dijkstra-Naive}(\mathbf{G}, w, s)$}
            \li $\proc{Init-Single-Source}(\mathbf{G}, s)$
            \li $\mathbf{S} \gets \emptyset$
            \li $\mathbf{Q} \gets \mathbf{V}$
            \li \While $\mathbf{Q} \neq \emptyset$
            \li   \Do $u \gets \proc{Extract-Min-Array}(\mathbf{Q})$
            \li       $\mathbf{S} \gets \mathbf{S} \cup \{u\}$
            \li       \For each vertice $v \in \mathbf{G}.Adj[u]$
            \li           \Do $\proc{Relax}(u, v, w)$
                    \End
                \End
            \end{codebox}
        \end{column}
    \end{columns}
\end{frame}


% Slide: Dijkstra con Priority Queue
\begin{frame}[fragile]{Miglioramento: \textsc{Dijkstra} con Priority Queue}
    \begin{columns}
        \begin{column}{0.5\textwidth}
            L'uso di un Min-Heap (o Coda di Priorità) velocizza l'estrazione, ma introduce un costo logaritmico fisso per il \textit{sorting} della frontiera.
            \vspace{0.2cm}
            \begin{block}{Complessità}
                \begin{itemize}
                    \item \textbf{Estrazione minimo:} $\mathcal{O}(\log \vert \mathbf{V}\vert)$
                    \item \textbf{Decrease Key:} $\mathcal{O}(\log \vert \mathbf{V}\vert)$ 
                    \item \textbf{Complessità Computazionale:} $\mathcal{O}(\vert \mathbf{E}\vert \log \vert \mathbf{V}\vert + \vert \mathbf{V}\vert \log \vert \mathbf{V}\vert)$ 
                \end{itemize}
            \end{block}
            \small \textit{(Su grafi sparsi $\vert \mathbf{E}\vert$ $\approx$ $\vert \mathbf{V}\vert$, il tempo è limitato a $\Omega(\vert \mathbf{V}\vert \log \vert \mathbf{V}\vert)$).}
        \end{column}
        \begin{column}{0.5\textwidth}
            \scriptsize
            \begin{codebox}
            \Procname{$\proc{Dijkstra-PQ}(\mathbf{G}, w, s)$}
            \li $\proc{Init-Single-Source}(\mathbf{G}, s)$
            \li $\mathbf{S} \gets \emptyset$
            \li $\mathbf{Q} \gets \emptyset$ \Comment Min-Heap
            \li \For each $v \in \mathbf{V}$ 
            \li \Do $\proc{Insert}$ $(\mathbf{Q}, v)$ \End
            \li \While $\mathbf{Q} \neq \emptyset$
            \li   \Do $u \gets \proc{Extract-Min}(\mathbf{Q})$
            \li       $\mathbf{S} \gets \mathbf{S} \cup \{u\}$
            \li       \For each $v \in \mathbf{G}.Adj[u]$
            \li           \Do $\proc{Relax-and-Decrease-Key}(u, v, w)$
                    \End
                \End
            \end{codebox}
        \end{column}
    \end{columns}
\end{frame}

% Slide: Operazione di rilassamento
\begin{frame}{Operazione di Rilassamento (Relax)}
    Sia il classico Dijkstra che il moderno BMSSP condividono la stessa operazione fondamentale per aggiornare le distanze: il \textbf{Rilassamento} di un arco.
    \vspace{0.2cm}

    \begin{block}{L'Idea Intuitiva}
        Immaginiamo di trovarci nel nodo $u$ e di voler andare nel nodo $v$ tramite l'arco $(u,v)$. \\
        Ci chiediamo: \textbf{"Passare da $u$ mi permette di arrivare a $v$ facendo meno strada?"}
    \end{block}

    \begin{exampleblock}{La Formula}
        \centering
        \textbf{se} $\text{dist}(u) + w(u,v) < \text{dist}(v)$ \textbf{allora} $\text{dist}(v) \gets \text{dist}(u) + w(u,v)$
    \end{exampleblock}

    \small
    \textbf{Legenda dei termini:}
    \begin{itemize}
        \item \textbf{$\text{dist}(v)$}: La miglior distanza trovata finora dalla sorgente al nodo $v$.
        \item \textbf{$w(u,v)$}: Il costo (peso) dell'arco diretto che collega il nodo $u$ al nodo $v$.
        \item \textbf{$\gets$}: Operazione di aggiornamento (assegnamento del nuovo valore).
    \end{itemize}
\end{frame}

% Slide: Collo di Bottiglia
\begin{frame}{Il Collo di Bottiglia: La Barriera del Sorting}
    \begin{block}{Il Limite di Dijkstra}
        \begin{itemize}
            \item L'algoritmo di Dijkstra mantiene una "frontiera" di vertici scoperti ma non ancora processati.
            \item Ad ogni passo, deve estrarre il vertice con la distanza minima usando una coda di priorità (Min-Heap).
            \item Questo continuo ordinamento richiede un tempo minimo di $\Omega(\vert \mathbf{V}\vert \log \vert \mathbf{V}\vert)$, creando un "collo di bottiglia" invalicabile sui grafi sparsi.
        \end{itemize}
    \end{block}
    \begin{block}{L'Idea del Paper}
        Fondere Dijkstra con l'approccio di \textbf{Bellman-Ford}: esplorare i vertici in "gruppi" delimitati da un upper-bound $B$, evitando di mantenere un ordine totale tra tutti i nodi scoperti.
    \end{block}
\end{frame}

\begin{frame}{Confronto 3 approcci}
    \begin{table}
        \centering
        \renewcommand{\arraystretch}{1.8}
        \begin{tabular}{|p{2.6cm}|p{3.2cm}|p{3.6cm}|p{3.2cm}|} 
        \hline
        \textbf{\textcolor{darkgreen}{Caratteristica}} & \textbf{\textcolor{darkgreen}{Dijkstra Ingenuo}} & \textbf{\textcolor{darkgreen}{Dijkstra (Min-Heap)}} & \textbf{\textcolor{darkgreen}{BMSSP}} \\ \hline
        \textbf{Complessità Temporale} & $\mathcal{O}(\vert \mathbf{V}\vert^2)$ & $\mathcal{O}(\vert \mathbf{E}\vert + \vert \mathbf{V}\vert \log \vert \mathbf{V}\vert)$ & $\mathcal{O}(\vert \mathbf{E}\vert \log^{2/3} \vert \mathbf{V}\vert)$ \\ \hline
        \textbf{Complessità Spaziale} & $\mathcal{O}(\vert \mathbf{V}\vert + \vert \mathbf{E}\vert)$ & $\mathcal{O}(\vert \mathbf{V}\vert + \vert \mathbf{E}\vert)$ & $\mathcal{O}(\vert \mathbf{V}\vert + \vert \mathbf{E}\vert)$ \\ \hline
        \end{tabular}
    \end{table}  
\end{frame}

% Slide: Assunzioni del paper
\begin{frame}{Introduzione al BMSSP: Assunzioni}
    \begin{block}{Assunzioni}
        \begin{itemize}
            \item ogni vertice $v$ è raggiungibile da $s$ (cioè $\vert \mathbf{E}\vert \geq \vert \mathbf{V}\vert-1$)
            \begin{itemize}
                \item \textcolor{red}{Se violata}: non succede niente, ma quest'assuzione serve matematicamente per semplificare le notazioni asintotiche.
            \end{itemize}
            \item Modello \textit{Comparison-Addition}: L'algoritmo lavora sui pesi degli archi usando solo operazioni di confronto e addizione.
\begin{itemize}
    \item \textcolor{red}{\textbf{Se violata:}} Se il modello permettesse manipolazioni di bit (modello \textit{Word RAM}) e i pesi fossero limitati ai \textbf{numeri interi}, questo algoritmo non sarebbe necessario. L'algoritmo di Thorup risolve già il problema in tempo quasi-lineare sfruttando le operazioni binarie. Il BMSSP è fondamentale proprio perché garantisce questi tempi sub-ottimali operando nel modello classico con \textbf{pesi reali} ($\mathbb{R}$).
\end{itemize}
            \item tutte le lunghezze dei percorsi sono diverse
\begin{itemize}
    \item \textcolor{red}{Se violata}: è un'assunzione senza perdita di generalità, per evitare la formazione dei DAG (Grafo Aciclico Diretto)
\end{itemize}

        \end{itemize}
    \end{block}
\end{frame}

% Slide: Algoritmo di Thorup
\begin{frame}{Digressione: Il funzionamento dell'Algoritmo di Thorup}
    Per contestualizzare il limite del BMSSP, è utile ricordare come opera l'algoritmo di Thorup, che risolve il SSSP in tempo quasi-lineare per grafi diretti.
    \vspace{0.1cm}
    
    \begin{block}{L'Idea Chiave: Bucket Sorting su bit}
        \small
        \begin{itemize}
            \setlength{\itemsep}{0.1cm} 
            \item \textbf{Struttura:} Invece del Min-Heap, distribuisce i vertici in \textbf{bucket} (secchi) basandosi sui bit più significativi delle distanze.
            \item \textbf{Operazioni Word RAM:} Utilizza operazioni a livello di bit (\textit{bitwise shift}, identificazione del bit più significativo) che la CPU esegue in tempo costante $\mathcal{O}(1)$.
            \item \textbf{Complessità:} Raggiunge tempi sub-ottimali ($\mathcal{O}(|E| + |V| \log \log |V|)$) abbattendo il costo di estrazione del minimo.
        \end{itemize}
    \end{block}
    
\end{frame}

% Slide: Thorup e l'Operatività sui Bit
\begin{frame}{Thorup: Operatività sui Bit}
    È fondamentale capire \textbf{come} Thorup usa i bit per aggirare l'ordinamento in tempo $\mathcal{O}(1)$.
    \vspace{0.1cm}
    
    \begin{block}{Calcolo dell'Indice del Bucket (XOR e MSB)}
        \small
        \begin{itemize}
            \setlength{\itemsep}{0.1cm} 
            \item Invece di un Min-Heap, i vertici sono smistati in \textbf{bucket} indicizzati.
            \item Quando l'arco $(u,v)$ viene rilassato, in quale bucket va spostato il nodo $v$?
            \item \textbf{La soluzione:} L'indice si calcola trovando il bit più significativo (\textit{MSB}) in cui la vecchia e la nuova distanza differiscono. \\
            Operativamente: \textbf{$\text{Indice} = \text{MSB}(\mathbf{d}_{\text{old}}[v] \oplus \mathbf{d}_{\text{new}}[v])$}
            \item L'operatore logico \textit{XOR} ($\oplus$) e la ricerca dell'MSB sono eseguiti in un singolo ciclo di clock della CPU nel modello \textit{Word RAM}.
        \end{itemize}
    \end{block}
    
    \vspace{0.1cm}
    \footnotesize \textbf{Il limite:} Questo trucco esige che i pesi siano \textbf{numeri interi} (sequenze fisse di bit). Il BMSSP usa pesi \textbf{reali} ($\mathbb{R}$), sui quali operatori a livello di bit come l'XOR non hanno senso o non sono applicabili.
\end{frame}

% Slide: Thorup: La Gerarchia dei Componenti
\begin{frame}{Thorup: La Gerarchia dei Componenti}
    Oltre ai bit, il secondo pilastro di Thorup è la gestione gerarchica del grafo per evitare l'ordinamento globale dei nodi.
    \vspace{0.3cm}

    \begin{columns}[T]
        \begin{column}{0.52\textwidth}
            \begin{block}{L'approccio Hierarchical}
                \footnotesize
                \begin{itemize}
                    \setlength{\itemsep}{0.1cm}
                    \item \textbf{Decomposizione:} Il grafo viene suddiviso in componenti in base ai pesi degli archi.
                    \item \textbf{Component Tree:} Si crea un albero: le foglie sono i nodi, i rami sono i gruppi.
                    \item \textbf{Visita Topologica:} Si visita l'albero in ordine topologico, non con un Min-Heap.
                \end{itemize}
            \end{block}
            \vspace{0.2cm}
            \footnotesize \textbf{Risultato:} Si estraggono "lotti" di nodi già pronti per il rilassamento, riducendo i confronti.
        \end{column}
        
         \begin{column}{0.44\textwidth}
        
            \centering
            \vspace{0.5cm}
            
            \resizebox{\linewidth}{!}{% 
            \begin{tikzpicture}[
                node distance=0.8cm,
                comp/.style={rectangle, draw=darkgreen, fill=darkgreen!10, thick, rounded corners, font=\scriptsize, minimum width=1.5cm},
                leaf/.style={circle, draw=gray, fill=gray!20, font=\tiny, inner sep=1pt}
            ]
                \node[comp] (root) {Radice (Tutto il Grafo)};
                \node[comp, below left=of root, xshift=-0.2cm] (c1) {Comp. $2^i$};
                \node[comp, below right=of root, xshift=0.2cm] (c2) {Comp. $2^j$};
                
                \node[leaf, below=0.4cm of c1, xshift=-0.4cm] (n1) {$v_1$};
                \node[leaf, below=0.4cm of c1, xshift=0cm] (n2) {$v_2$};
                \node[leaf, below=0.4cm of c1, xshift=0.4cm] (n3) {$v_3$};
                
                \draw[->, thick, gray] (root) -- (c1);
                \draw[->, thick, gray] (root) -- (c2);
                \draw[-] (c1) -- (n1);
                \draw[-] (c1) -- (n2);
                \draw[-] (c1) -- (n3);
                
                \node[below=0.1cm of n2, font=\tiny\itshape] {Nodi individuali};
            \end{tikzpicture}%
            }
        \end{column}
    \end{columns}
\end{frame}

% Slide: Thorup Running Example - Setup
\begin{frame}{Thorup Running Example: 1. Lo scenario}
    Consideriamo un nodo $v$ con una distanza provvisoria dalla sorgente.
    \vspace{0.3cm}
    \begin{columns}[T]
        \begin{column}{0.5\textwidth}
            \begin{itemize}
                \item \textbf{Distanza corrente:} $d_{old}(v) = 13$ (binario: \texttt{1101})
                \item \textbf{Nuovo arco:} Arriva un arco che abbassa la distanza a $d_{new}(v) = 9$ (binario: \texttt{1001})
                \item \textbf{Il Problema:} In Dijkstra dovremmo riordinare la coda. Thorup deve solo decidere in quale "secchio" (\textit{bucket}) spostare $v$.
            \end{itemize}
        \end{column}
        \begin{column}{0.45\textwidth}
            \centering
            % Qui puoi inserire un piccolo schema o testo riassuntivo
            \begin{block}{Modello Word RAM}
                La CPU vede i pesi come stringhe di bit, permettendo calcoli logici istantanei che ignorano il valore decimale.
            \end{block}
    \end{column}
    \end{columns}
\end{frame}

% Slide: Thorup Running Example Calcolo
\begin{frame}{Thorup Running Example: 2. Il calcolo del Bucket}
    Per trovare il bucket corretto senza confronti, Thorup calcola la differenza bit a bit.
    \vspace{0.3cm}
    \begin{exampleblock}{Calcolo XOR e MSB}
        \centering
        \begin{tabular}{rcc}
            $d_{old}(v) = 13$ & $\rightarrow$ & \texttt{1 1 0 1} \\
            $d_{new}(v) = 9$  & $\rightarrow$ & \texttt{1 0 0 1} \\ \hline
            \textbf{XOR ($\oplus$)} & $\rightarrow$ & \texttt{0 \underline{1} 0 0} \\
        \end{tabular}
    \end{exampleblock}
    \vspace{0.3cm}
    \begin{itemize}
        \item \textbf{Risultato:} Il primo bit che differisce è il secondo da sinistra (valore posizionale $2^2$).
        \item \textbf{Azione:} Il nodo $v$ viene spostato nel bucket di livello $2$.
        \item \textbf{Efficienza:} Questa operazione è eseguita dalla CPU in un istante ($\mathcal{O}(1)$), indipendentemente da quanti altri nodi ci siano nel sistema.
    \end{itemize}
\end{frame}

% Slide: Confronto con Thorup
\begin{frame}{Confronto tra BMSSP e l'Algoritmo di Thorup}
    Perché il BMSSP è rivoluzionario se Thorup (1999) aveva già risolto il SSSP in tempo quasi-lineare? La differenza risiede nel \textbf{modello computazionale}.
   \vspace{-0.3cm}
    \begin{table}[]
        \centering
        \renewcommand{\arraystretch}{1.5}
        \begin{tabular}{|p{3.8cm}|p{4.5cm}|p{4.5cm}|}
        \hline
        \textbf{\textcolor{darkgreen}{Caratteristica}} & \textbf{\textcolor{darkgreen}{Algoritmo di Thorup}} & \textbf{\textcolor{darkgreen}{Algoritmo BMSSP}} \\ \hline
        \textbf{Modello} & Word RAM & Comparison-Addition \\ \hline
        \textbf{Dominio dei Pesi} & Numeri \textbf{Interi} ($\mathbb{N}$) & Numeri \textbf{Reali} ($\mathbb{R}_{\ge 0}$) \\ \hline
        \textbf{Operazioni Ammesse} & Somme, Confronti, \textbf{Shift di bit, operatori logici} & Esclusivamente somme e confronti ($+$, $<$) \\ \hline
        \textbf{Logica di estrazione} & \textit{Bucket Sort} sui prefissi binari & \textit{Batch Prepend} su Red-Black Tree \\ \hline
        \end{tabular}
    \end{table}
        Thorup accelera l'esplorazione guardando la rappresentazione binaria dei pesi interi in memoria. Il BMSSP, invece, tratta i pesi reali come "scatole chiuse".
\end{frame}

% Slide: Notazioni
\begin{frame}{Introduzione al BMSSP: Notazioni}
    \begin{block}{Notazioni}
    \begin{itemize}
        \item $d(u)$: lunghezza del shortest path da $s$ a $u$
        \item $\hat{d}[u]$: variabile globale per stimare $d(u)$, che viene inizializzata $\hat{d}[s]=0$ e $\hat{d}[v]=\infty$ per i vertici che non sono la sorgente $s$
        \item $u$ è detto complete: se $\hat{d}[u] = d(u)$ (altrimenti è detto incomplete)
        \item $S$ è un set complete: se tutti i suoi vertici lo sono.
        \item $Pred[v]$: predecessore di $v$
        \item $\langle l,\alpha,v_\alpha,v_{\alpha-1},...,v_1 \rangle$: percorso di lunghezza $l$, che collega $v_1 = s$ con $v_\alpha$ attraversando $\alpha$ vertici
    \end{itemize}
    \end{block}
\end{frame}

% Slide: L'Intuizione: Ridurre la Frontiera
\begin{frame}{L'Intuizione: Ridurre la Frontiera}
    Invece di ordinare una frontiera enorme di grandezza $\Theta(\vert \mathbf{V}\vert)$, l'algoritmo restringe la ricerca a un piccolo set di \textbf{Pivot}.
    
    \vspace{0.5cm}
    \begin{center}
    \begin{tikzpicture}[
        node distance=1cm,
        box/.style={rectangle, rounded corners, draw=darkgreen, fill=blockbodybg, thick, minimum width=2.5cm, minimum height=1cm, align=center},
        arrow/.style={-{Stealth[scale=1.2]}, thick, darkgreen}
        ]
        
        \node[box] (frontiera) {Frontiera Gigante\\ $\Theta(\vert \mathbf{V}\vert)$ nodi};
        \node[box, right=of frontiera, fill=lightgreen!30] (bellman) {Rilassamento\\ ($k$ step di Bellman-Ford)};
        \node[box, right=of bellman, draw=red!80, fill=red!10] (pivot) {Insieme di Pivot\\ $\le |U|/k$ nodi};
        
        \draw[arrow] (frontiera) -- (bellman);
        \draw[arrow] (bellman) -- (pivot);
        
    \end{tikzpicture}
    \end{center}
    \vspace{0.1cm}
    
    \begin{itemize}
        \item Se un percorso richiede meno di $k$ step, viene risolto subito con Bellman-Ford.
        \item I nodi rimanenti dipendono da pochi \textbf{Pivot} (le radici degli alberi dei cammini minimi di profondità $\ge k$).
    \end{itemize}
    \begin{block}{Rilassamento}
        \centering
        \textbf{se} $\mathbf{d}[u] + w_{uv} < \mathbf{d}[v]$ \textbf{allora} $\mathbf{d}[v] \gets \mathbf{d}[u] + w_{uv}$
    \end{block}
\end{frame}

% Slide: Algoritmo BMSSP
\begin{frame}{Bounded Multi-Source Shortest Path (BMSSP)}
    Il cuore dell'innovazione è la procedura ricorsiva \textbf{BMSSP(l, B, S)}:\\
    Introduciamo prima il concetto di livello $l$:\begin{itemize}
            \item \textbf{Cos'è $l$?} È l'\textbf{indice del livello di ricorsione}. \\Parte da un valore massimo e decresce unitariamente ad ogni chiamata interna.
            \item \textbf{Controllo della Granularità:} Il livello $l$ determina l'ampiezza massima dei blocchi ($M = 2^{(l-1)t}$) all'interno della struttura dati $\mathbf{D}$.
            \begin{itemize}
                \item \textbf{Alti livelli di $l$:} Blocchi molto grandi. Bassa precisione di ordinamento, ma massima velocità per l'inserimento a lotti (\textit{Batch Prepend}).
                \item \textbf{Bassi livelli di $l$:} Blocchi sempre più piccoli e precisi.
            \end{itemize}
            \item \textbf{Caso Base ($l=0$):} L'ampiezza del blocco collassa. La ricorsione termina e l'algoritmo ripiega su un mini-Dijkstra esatto limitato a pochi vertici.
        \end{itemize}
\end{frame}

% Slide: Algoritmo BMSSP
\begin{frame}{Bounded Multi-Source Shortest Path (BMSSP)}
    \begin{block}{Divide et Impera}
        \begin{itemize}
            \item \textbf{Input:} un livello $l$, un upper bound $B$, e un set di sorgenti $S$.
            \item \textbf{Logica:} L'algoritmo partiziona adattivamente il problema in sotto-problemi di dimensioni inferiori.
            \item \textbf{Ricorsione:} Scende fino al caso base ($l=0$) dove esegue un "mini-Dijkstra" limitato all'esplorazione di $k+1$ nodi.
            \item \textbf{Output:} Trova tutte le distanze reali inferiori a un nuovo bound $B' \le B$.
        \end{itemize}
    \end{block}
\end{frame}

% Slide: La Struttura Dati Ibrida
\begin{frame}{La Struttura Dati Ibrida (Lemma 3.3)}
    Per gestire i partizionamenti in modo efficiente, il paper introduce una struttura dati ad-hoc $\mathbf{D}$ composta da:
    \vspace{0.3cm}
    
    \begin{itemize}
        \item \textbf{Blocchi (Liste) $\mathbf{D}_l$: } I nodi con distanze simili vengono raggruppati in blocchi di dimensione massima $M$. (chiamata $\mathbf{D}_{0}$ nel paper)
        \item \textbf{Albero Red-Black $\mathbf{D}_{rb}$: } Mantiene i limiti superiori (\textit{upper bounds}) di ciascun blocco. (chiamata $\mathbf{D}_{1}$ nel paper)
        \item \textbf{Inserimento a Lotti (\textit{Batch Prepend}):} Invece di inserire e ordinare i nuovi nodi uno ad uno (come fa la coda di priorità in Dijkstra), i nodi con distanze simili vengono raggruppati in un unico blocco. Questo intero blocco viene inserito tutto in una volta in testa alla struttura, aggirando il costo del sorting continuo.
    \end{itemize}
\end{frame}

% Slide: Esempio Dijkstra
\begin{frame}{Esempio Visivo: L'approccio Classico (Dijkstra)}
    Dijkstra esplora il grafo rigorosamente in base alla distanza minima.
    \vspace{0.1cm}
    
    \begin{center}
    \begin{tikzpicture}[
        node distance=1cm and 1.5cm,
        vertex/.style={circle, draw=darkgreen, fill=blockbodybg, thick, minimum size=0.6cm, inner sep=1pt},
        edge/.style={-{Stealth}, thick},
        pq/.style={rectangle, draw=red, fill=red!10, thick, rounded corners, align=center, font=\footnotesize}
        ]
        
        % Vertici
        \node[vertex, fill=darkgreen, text=white] (S) {$s$};
        \node[vertex, above right=of S] (A) {$A$};
        \node[vertex, below right=of S] (B) {$B$};
        \node[vertex, right=of A] (C) {$C$};
        \node[vertex, right=of B] (D) {$D$};
        
        % Archi con pesi
        \draw[edge] (S) -- node[above left, scale=0.8] {2} (A);
        \draw[edge] (S) -- node[below left, scale=0.8] {5} (B);
        \draw[edge] (A) -- node[above, scale=0.8] {1} (C);
        \draw[edge] (B) -- node[below, scale=0.8] {1} (D);
        \draw[edge] (A) -- node[left, scale=0.8] {4} (B);
        
        % Box Priority Queue
        \node[pq, right=1.2cm of C, yshift=-0.8cm] (PQ) {
            \textbf{Priority Queue}\\ 
            \textit{Costo: $\mathcal{O}(\log n)$}\\
            \rule{2.5cm}{0.4pt}\\
            1. Estrarre $A$ (d=2)\\
            2. Inserire $C$ (d=3)\\
            3. Ordinare $B$ e $C$\\
            ...
        };
        
        % Frecce
        \draw[->, red, thick, dashed] (A) edge[bend left] (PQ);
        \draw[->, red, thick, dashed] (B) edge[bend right] (PQ);
        
    \end{tikzpicture}
    \end{center}
    
    \textbf{Problema:} Anche se l'arco $A \to C$ è ovvio, Dijkstra deve \textbf{sempre} rimettere tutto in coda e riordinare l'intera frontiera.
\end{frame}

% Slide: Esempio Nuovo Algoritmo
\begin{frame}{Esempio Visivo: Il Nuovo Approccio (BMSSP)}
    Si usano i $k$ step : raggruppa l'esplorazione ed evita il sorting globale.
    \begin{center}
    \begin{tikzpicture}[
        node distance=1cm and 1.5cm,
        vertex/.style={circle, draw=darkgreen, fill=blockbodybg, thick, minimum size=0.6cm, inner sep=1pt},
        edge/.style={-{Stealth}, thick}
        ]
        
        % Vertici
        \node[vertex, fill=darkgreen, text=white] (S) {$s$};
        \node[vertex, above right=of S] (A) {$A$};
        \node[vertex, below right=of S] (B) {$B$};
        \node[vertex, right=of A] (C) {$C$};
        \node[vertex, right=of B] (D) {$D$};
        
        % Archi con pesi
        \draw[edge] (S) -- node[above left, scale=0.8] {2} (A);
        \draw[edge] (S) -- node[below left, scale=0.8] {5} (B);
        \draw[edge] (A) -- node[above, scale=0.8] {1} (C);
        \draw[edge] (B) -- node[below, scale=0.8] {1} (D);
        \draw[edge] (A) -- node[left, scale=0.8] {4} (B);

        \draw[red, thick] (A) circle (0.5cm);
        \node[red, font=\footnotesize] at ([yshift=-0.7cm]A) {\textbf{Pivot}};
        
        \draw[blue, dashed, thick, rounded corners] ([xshift=-0.7cm,yshift=0.7cm]A) rectangle ([xshift=0.7cm,yshift=-0.7cm]C);
        \node[blue, font=\footnotesize] at ([yshift=0.9cm, xshift=0.8cm]A) {Relax per $k$ step};
        
    \end{tikzpicture}
    \end{center}
    \begin{itemize}
        \item \textbf{Propagazione senza Priority Queue:} Si \textbf{rilassano} gli archi a cascata per $k$ passi. 
        \item Il nodo $A$ genera un sotto-albero abbastanza profondo da essere promosso a \textbf{Pivot}. Solo i Pivot passano alla ricorsione successiva.
    \end{itemize}
\end{frame}





% Slide: Pseudocodice BMSSP
\begin{frame}[fragile]{Pseudocodice del Bounded Multi-Source (\textsc{Bmssp})}
    \begin{columns}
        \begin{column}{0.45\textwidth}
            La procedura partiziona l'esplorazione ricorsivamente.
            \vspace{0.2cm}
            \begin{itemize}
                \item Usa $\textsc{FindPivots}$ per isolare $\mathbf{P}$.
                \item Estrae i lotti di minimi con $\mathbf{D}.\textsc{Pull}()$.
                \item Scende di livello fino a ripiegare su un mini-Dijkstra limitato.
            \end{itemize}
        \end{column}
        \begin{column}{0.55\textwidth}
                \vspace{-0.3cm}

            \scriptsize
            \begin{codebox}
            \Procname{$\proc{Bmssp}(l, B, \mathbf{S})$}
            \li \If $l = 0$ 
            \li \Then \Return $\proc{BaseCase}(B, \mathbf{S})$ \End
            \li $\mathbf{P}, \mathbf{W} \gets \proc{FindPivots}(B, \mathbf{S})$
            \li $\mathbf{D}.\proc{Init}(M, B)$
            \li \For each $x \in \mathbf{P}$ 
            \li \Do $\mathbf{D}.\proc{Insert}(\langle x, \mathbf{d}[x] \rangle)$ \End
            \li \While $|\mathbf{U}| < k 2^{lt}$ e $\mathbf{D} \neq \emptyset$ \Do
            \li   $B_i, \mathbf{S}_i \gets \mathbf{D}.\proc{Pull}()$
            \li   $B'_i, \mathbf{U}_i \gets \proc{Bmssp}(l-1, B_i, \mathbf{S}_i)$
            \li   $\mathbf{U} \gets \mathbf{U} \cup \mathbf{U}_i$
            \li   $\mathbf{K} \gets \emptyset$
            \li   \For each $(u,v)$ uscente da $\mathbf{U}_i$ \Do
            \li       \If $\proc{Relax}(u, v)$ \Then
            \li           \If $\mathbf{d}[v] \in [B_i, B)$ 
            \li \Then $\mathbf{D}.\proc{Insert}(v)$
            \li           \ElseIf $\mathbf{d}[v] \in [B'_i, B_i)$ 
            \li \Then $\mathbf{K} \gets \mathbf{K} \cup \{v\}$
                          \End
                      \End
                  \End
            \li   $\mathbf{D}.\proc{BatchPrepend}(\mathbf{K})$
                \End
            \li \Return $B', \mathbf{U}$
            \end{codebox}
        \end{column}
    \end{columns}
\end{frame}


% Slide: Running Example 1
\begin{frame}{Running Example: 1. Fase di $\textsc{FindPivots}$}
    Fissiamo la soglia $k=2$. Iniziamo con il limite superiore $B = \infty$ e la sorgente $\mathbf{S} = \{s\}$. Il primo passo è propagare i cammini minimi provvisori per $k$ iterazioni.
    \vspace{0.3cm}
    
    \begin{columns}
        \begin{column}{0.45\textwidth}
            \begin{itemize}
                \item \textbf{Step 1:} Rilassiamo da $s$. Troviamo $A$ ($\mathbf{d}[A]=2$) e $B$ ($\mathbf{d}[B]=5$).
                \item \textbf{Step 2:} Rilassiamo da $A$ e $B$. Troviamo $C$ ($\mathbf{d}[C]=3$) ed $E$ ($\mathbf{d}[E]=7$).
                \item \textbf{Risultato:} L'insieme dei vertici scoperti è $\mathbf{W} = \{s, A, B, C, E\}$.
                \item \textbf{Nota:} Nessun ordinamento totale è stato effettuato.
            \end{itemize}
        \end{column}
        \begin{column}{0.55\textwidth}
            \begin{center}
            \begin{tikzpicture}[
                node distance=1.5cm and 1.5cm,
                vertex/.style={circle, draw=darkgreen, fill=blockbodybg, thick, minimum size=0.6cm, inner sep=1pt},
                edge/.style={-{Stealth}, thick},
                relaxed/.style={-{Stealth}, thick, blue}
                ]
                
                \node[vertex, fill=darkgreen, text=white] (S) {$s$};
                \node[vertex, above right=of S] (A) {$A$};
                \node[vertex, right=of A] (C) {$C$};
                \node[vertex, right=of C] (D) {$D$};
                \node[vertex, below right=of S] (B) {$B$};
                \node[vertex, right=of B] (E) {$E$};
                
                % Archi Rilassati (Blu)
                \draw[relaxed] (S) -- node[above left, scale=0.8] {2} (A);
                \draw[relaxed] (S) -- node[below left, scale=0.8] {5} (B);
                \draw[relaxed] (A) -- node[above, scale=0.8] {1} (C);
                \draw[relaxed] (B) -- node[below, scale=0.8] {2} (E);
                
                % Archi non ancora esplorati (Neri)
                \draw[edge] (C) -- node[above, scale=0.8] {1} (D);
                \draw[edge] (A) -- node[left, scale=0.8] {6} (B);
            \end{tikzpicture}
            \end{center}
        \end{column}
    \end{columns}
\end{frame}

% Slide: Running Example 2
\begin{frame}{Running Example: 2. Foresta $\mathbf{F}$ e Selezione Pivot}
    Costruiamo la foresta $\mathbf{F}$ analizzando i predecessori ($\mathbf{pred}$) appena scoperti in $\mathbf{W}$.
    \vspace{0.3cm}
    
    \begin{columns}
        \begin{column}{0.45\textwidth}
            \begin{itemize}
                \item L'albero radicato nella sorgente $s$ ha dimensione $5$.
                \item Poiché la dimensione è $\ge k$ ($5 \ge 2$), il vertice $s$ viene promosso a \textbf{Pivot}.
                \item Insieme dei Pivot: $\mathbf{P} = \{s\}$.
                \item Se un nodo genera un albero piccolo ($< k$), non diventa Pivot e viene rimandato alla ricorsione.
            \end{itemize}
        \end{column}
        \begin{column}{0.55\textwidth}
            \begin{center}
            \begin{tikzpicture}[
                node distance=1.5cm and 1.5cm,
                vertex/.style={circle, draw=darkgreen, fill=blockbodybg, thick, minimum size=0.6cm, inner sep=1pt},
                pivot/.style={circle, draw=red, fill=red!20, thick, minimum size=0.6cm, inner sep=1pt},
                tree/.style={-{Stealth}, ultra thick, darkgreen}
                ]
                
                \node[pivot] (S) {$s$};
                \node[vertex, above right=of S] (A) {$A$};
                \node[vertex, right=of A] (C) {$C$};
                \node[vertex, right=of C, draw=gray, text=gray] (D) {$D$}; 
                \node[vertex, below right=of S] (B) {$B$};
                \node[vertex, right=of B] (E) {$E$};
                
                \node[red, font=\footnotesize] at ([yshift=-0.6cm]S) {\textbf{Pivot}};
                
                % Foresta F
                \draw[tree] (S) -- (A);
                \draw[tree] (S) -- (B);
                \draw[tree] (A) -- (C);
                \draw[tree] (B) -- (E);
                
                \draw[-{Stealth}, gray, dashed] (C) -- (D);
                \draw[-{Stealth}, gray, dashed] (A) -- (B);
            \end{tikzpicture}
            \end{center}
        \end{column}
    \end{columns}
\end{frame}

% Slide: Running Example 3
\begin{frame}{Running Example: 3. Setup di $\mathbf{D}$ e Prima \textsc{Pull}}
    L'algoritmo istanzia la struttura dati ibrida $\mathbf{D}$ e vi inserisce i Pivot per prepararli alla ricorsione.
    \vspace{0.3cm}
    
    \begin{columns}
        \begin{column}{0.5\textwidth}
            \begin{itemize}
                \item $\mathbf{D}.\textsc{Init}(\infty)$: Si crea un RB-Tree per $\mathbf{D}_{rb}$ e una lista vuota per $\mathbf{D}_l$.
                \item $\mathbf{D}.\textsc{Insert}(\langle s, 0 \rangle)$ inserisce la sorgente nell'albero $\mathbf{D}_{rb}$.
                \item \textbf{Inizio Ciclo While:} 
                \item Eseguiamo $\mathbf{D}.\textsc{Pull}()$: estraiamo i minimi.
                \item \textbf{Risultato:} $\mathbf{S}_1 = \{s\}$. Poiché $\mathbf{D}$ ora è vuota, il nuovo limite temporaneo è $B_1 = \infty$.
            \end{itemize}
        \end{column}
        \begin{column}{0.5\textwidth}
            \begin{center}
            \begin{tikzpicture}[
                node distance=0.8cm,
                block/.style={rectangle, draw=darkgreen, fill=blockbodybg, thick, rounded corners, minimum width=2cm, minimum height=0.6cm, align=center, font=\footnotesize},
                arrow/.style={-{Stealth}, thick}
                ]
                
                \node[block] (D1) {$\mathbf{D}_{rb}$ (RB-Tree)\\ $UB: \infty$};
                \node[block, right=0.5cm of D1, draw=blue!50] (D0) {$\mathbf{D}_l$ (Lista)\\ \textit{Vuota}};
                
                \node[block, below=0.5cm of D1, fill=red!10, draw=red] (Elem) {$\langle s, 0 \rangle$};
                \draw[arrow] (D1) -- (Elem);
                
                \node[rectangle, draw=red, thick, rounded corners, below=0.5cm of Elem, align=center, font=\footnotesize] (Pull) {
                    \textbf{\textsc{Pull}()}\\
                    $\mathbf{S}_1 = \{s\}$\\
                    $B_1 = \infty$
                };
                
                \draw[arrow, red, dashed, thick] (Elem) -- (Pull);
            \end{tikzpicture}
            \end{center}
        \end{column}
    \end{columns}
\end{frame}

% Slide: Running Example 4
\begin{frame}{Running Example: 4. Chiamata Ricorsiva e Ritorno}
    L'algoritmo invoca ricorsivamente $\textsc{Bmssp}(l-1, B_1, \mathbf{S}_1)$ per completare l'esplorazione attorno al pivot $s$.
    \vspace{0.3cm}
    
    \begin{columns}
        \begin{column}{0.45\textwidth}
            \begin{itemize}
                \item La ricorsione scende di livello e consolida i cammini minimi locali.
                \item \textbf{Ritorno della Ricorsione:}
                \item Restituisce i vertici completati $\mathbf{U}_1 = \{s, A, C\}$.
                \item Restituisce un nuovo limite $B'_1 = 4$ (le distanze in $\mathbf{U}_1$ sono $\le 4$).
            \end{itemize}
        \end{column}
        \begin{column}{0.55\textwidth}
            \begin{center}
            \begin{tikzpicture}[
                node distance=1.5cm and 1.5cm,
                vertex/.style={circle, draw=darkgreen, fill=blockbodybg, thick, minimum size=0.6cm, inner sep=1pt},
                completed/.style={circle, draw=purple, fill=purple!20, thick, minimum size=0.6cm, inner sep=1pt},
                edge/.style={-{Stealth}, thick}
                ]
                
                \node[completed] (S) {$s$};
                \node[completed, above right=of S, label={[purple]above:d=2}] (A) {$A$};
                \node[completed, right=of A, label={[purple]above:d=3}] (C) {$C$};
                \node[vertex, right=of C] (D) {$D$};
                \node[vertex, below right=of S] (B) {$B$};
                \node[vertex, right=of B] (E) {$E$};
                
                \draw[edge, purple, ultra thick] (S) -- (A);
                \draw[edge, purple, ultra thick] (A) -- (C);
                
                \node[purple, font=\footnotesize] at ([yshift=-1cm]A) {Insieme Completato $\mathbf{U}_1$};
                \node[purple, font=\footnotesize] at ([yshift=-1.4cm]A) {$B'_1 = 4$};
                
            \end{tikzpicture}
            \end{center}
        \end{column}
    \end{columns}
\end{frame}

% Slide: Running Example 5
\begin{frame}{Running Example: 5. Partizionamento dei Nuovi Nodi}
    Ora si esplorano gli archi \textbf{uscenti} dai nodi appena completati $\mathbf{U}_1$ e si classificano le nuove distanze in base a $B_1$ e $B'_1$.
    \vspace{0.3cm}
    
    \begin{columns}
        \begin{column}{0.55\textwidth}
            \begin{itemize}
                \item Arco $s \xrightarrow{5} B$: Troviamo $\mathbf{d}[B] = 5$.
                \item Arco $C \xrightarrow{1} D$: Troviamo $\mathbf{d}[D] = 4$.
                \item \textbf{Controllo dei Range:}
                \item Ricordiamo che $B'_1 = 4$ e $B_1 = \infty$.
                \item Entrambi i valori ($4$ e $5$) ricadono nel range $[\textcolor{red}{B'_1}, \textcolor{blue}{B_1}) = [\textcolor{red}{4}, \textcolor{blue}{\infty})$.
                \item \textbf{Regola del Paper:} Vengono parcheggiati nell'insieme temporaneo $\mathbf{K} = \{B, D\}$.
            \end{itemize}
        \end{column}
        \begin{column}{0.45\textwidth}
            \begin{center}
            \begin{tikzpicture}[
                node distance=1.5cm and 1.5cm,
                completed/.style={circle, draw=purple, fill=purple!20, thick, minimum size=0.6cm, inner sep=1pt},
                newnode/.style={circle, draw=orange, fill=orange!20, thick, minimum size=0.6cm, inner sep=1pt},
                edge/.style={-{Stealth}, thick}
                ]
                
                \node[completed] (S) {$s$};
                \node[completed, above right=of S] (A) {$A$};
                \node[completed, right=of A] (C) {$C$};
                \node[newnode, right=of C, label={[orange]right:d=4}] (D) {$D$};
                \node[newnode, below right=of S, label={[orange]right:d=5}] (B) {$B$};
                
                \draw[edge, purple] (S) -- (A);
                \draw[edge, purple] (A) -- (C);
                
                \draw[-{Stealth}, orange, ultra thick] (S) -- node[below left, text=black, scale=0.8] {5} (B);
                \draw[-{Stealth}, orange, ultra thick] (C) -- node[above, text=black, scale=0.8] {1} (D);
                
            \end{tikzpicture}
            \end{center}
        \end{column}
    \end{columns}
\end{frame}

% Slide: Running Example 6
\begin{frame}{Running Example: 6. $\textsc{BatchPrepend}$ e Loop}
    Per non dover fare $n$ ordinamenti in $\mathbf{D}_{rb}$, i vertici in $\mathbf{K}$ vengono convertiti in un blocco e inseriti in testa a $\mathbf{D}_l$ in tempo quasi $\mathcal{O}(1)$.
    \vspace{0.3cm}
    
    \begin{columns}
        \begin{column}{0.45\textwidth}
            \begin{itemize}
                \item Eseguiamo $\mathbf{D}.\textsc{BatchPrepend}(\mathbf{K})$.
                \item $\mathbf{D}_l$ riceve un nuovo blocco contenente $B$ e $D$.
                \item \textbf{Fine prima iterazione.} 
                \item Al giro successivo del ciclo \texttt{while}, la funzione $\textsc{Pull}()$ estrarrà il minimo proprio da $\mathbf{D}_l$ senza dover cercare nell'intero albero!
            \end{itemize}
        \end{column}
        \begin{column}{0.55\textwidth}
            \begin{center}
            \begin{tikzpicture}[
                node distance=0.8cm,
                block/.style={rectangle, draw=darkgreen, fill=blockbodybg, thick, rounded corners, minimum width=2.5cm, minimum height=0.6cm, align=center, font=\footnotesize},
                newgrp/.style={rectangle, draw=orange, fill=orange!10, thick, rounded corners, align=center, font=\footnotesize},
                arrow/.style={-{Stealth}, thick}
                ]
                
                \node[block] (D1) {$\mathbf{D}_{rb}$ (RB-Tree)\\ \textit{Vuoto}};
                \node[block, right=0.5cm of D1, draw=blue!50] (D0) {$\mathbf{D}_l$ (Lista)};
                
                \node[newgrp, below=0.5cm of D0] (K) {Blocco $\mathbf{K}$\\ $\langle D, 4 \rangle, \langle B, 5 \rangle$};
                
                \draw[arrow, orange, ultra thick] (K) -- node[right, scale=0.8] {\textsc{Prepend}} (D0);
                
            \end{tikzpicture}
            \end{center}
            \vspace{0.2cm}
            \begin{block}{Il Segreto della Complessità}
            Raggruppando i vertici in $\mathbf{K}$ e saltando la coda a priorità totale, si passa da $\Omega(n \log n)$ a $\mathcal{O}(m \log^{2/3} n)$.
            \end{block}
        \end{column}
    \end{columns}
\end{frame}










\end{document}
